\documentclass[aspectratio=169,usepdftitle=true,11pt,ngerman,t]{beamer}

\newcommand\lqfirstname{Name}
\newcommand\lqfullname{Name Last}
\newcommand\lqemail{name.last}
\errorcontextlines999999

% base
\usepackage[utf8]{inputenc}
\usepackage[T1]{fontenc}
\usepackage{lmodern}
\usepackage{microtype}
\usepackage[main=ngerman]{babel}
\usepackage{csquotes}

% graphics
\usepackage[prefix=]{xcolor-material}
\usepackage{graphics}
\usepackage{tcolorbox}
\usepackage{twemojis}
\usepackage{tikz}
\usetikzlibrary{
    positioning,
    calc,
    tikzmark,
    circuits.logic.IEC
}

% math
\usepackage{amsmath,amssymb}
\usepackage{mathtools}
\usepackage{array}
\usepackage{booktabs}
\usepackage{cancel}
\usepackage{ulem}

% flo
\usepackage{fancyqr}
\usepackage[addons]{color-palettes}
\usepackage[encoding,noloadlangs,cpalette,numinpar,fakeminted]{sopra-listings}
\usepackage{code-animation}
\usetheme[theorems]{colorful-dream}
\theoremstyle{exercise}
\newtheorem*{exercise*}{Aufgabe}
\theoremstyle{solve}
\newtheorem*{solve*}{Lösung}

% tools
\usepackage{algpseudocode}
\usepackage{dirtree}

% warnings
\usepackage{silence}
\WarningFilter{latexfont}{Some font shapes were not available}
\WarningFilter{latexfont}{Font shape}
\WarningFilter{pgf}{Returning node center instead of a point on node border}
\let\origenddocument\enddocument\renewcommand{\enddocument}{\hfuzz=10.95pt\origenddocument}
% title image: box
% ([vshift], inverse scale, width, content)
\newcommand\titleimagebox[4][0pt]{%
\newlength\tibheight\setlength\tibheight{4cm}
\newlength\tibtY\setlength\tibtY{-6em}
\titleimage{\begin{minipage}[c][#2\tibheight]{#3}\vspace*{#2\tibtY}%
    \vspace*{#1}\begin{tcolorbox}[boxrule=0.1pt]%
    \centering#4% TODO: bug: not centered if too wide
    \end{tcolorbox}%
\end{minipage}}}

% outro image: text
% (text)
\newcommand\outroimagetext[1]{\outro{Karlsruhe, \today%
\begin{tikzpicture}[remember picture,overlay]
    \node[align=center] at (current page.center) {#1};
\end{tikzpicture}}}% chktex 31

% reset counters
\newcommand\resetframecounters{\addtocounter{exercise}{-1}\addtocounter{solve}{-1}}
\newcommand\resetex{\addtocounter{exercise}{-1}}
\newcommand\resetsolve{\addtocounter{solve}{-1}}

% small footnote
\newcommand{\nicefootnote}[1]{\tikz[overlay,remember picture]{\node[anchor=south,text width=\dimexpr\textwidth+2em\relax,align=center,above=1em of current page.south] (nicefootnote) {\hrule\vspace{1pt}\footnotesize#1};}}

% kv
% ([draw], variables, values)
% {currently only works correctly for 4 variables}
\newcommand{\lqkv}[3][]{
    \begin{tikzpicture}
        \foreach \v [count=\k from 1] in {#2} {\global\let\Vdim=\k}
        \pgfmathsetmacro\rows{2^ceil(\Vdim/2)}
        \pgfmathsetmacro\cols{2^int(\Vdim/2)}

        \foreach \v [count=\k from 1] in {#2} {
            \pgfmathsetmacro\shift{0.5+floor(\k*ceil(\Vdim/2)/2)}

            \ifcase\k\relax
            \or
                \draw (\shift-1,0.75) -- (\shift+1,0.75);
                \node[fill=white] at (\shift,0.75) {\v};
            \or
                \draw (-0.75,-\shift) -- (-0.75,-\shift+2);
                \node[fill=white] at (-0.75,-\shift+1) {\v};
            \or
                \draw (\shift-2,-\cols+0.25) -- (\shift,-\cols+0.25);
                \node[fill=white] at (\shift-1,-\cols+0.25) {\v};
            \or
                \draw (\cols-0.25,-\shift+1) -- (\cols-0.25,-\shift+3);
                \node[fill=white] at (\cols-0.25,-\shift+2) {\v};
            \fi
        }
      
        \foreach \v [count=\k from 0] in {#3} {
            \pgfmathsetmacro\x{mod(\k,\rows)}
            \pgfmathsetmacro\y{int(\k/\rows)}
            \node[draw, minimum width=1cm, minimum height=1cm, align=center] at (\x,-\y) {\v};
        }

        \begin{scope}
            \clip (-0.5,0.5) rectangle (3.5,-3.5);
            \foreach \r in {#1} {\draw[rounded corners=1ex] \r;}
        \end{scope}

        \global\let\Vdim\undefined
    \end{tikzpicture}
}
\title{Digitaltechnik}
\institute{Karlsruher Institut für Technologie}

\author{\lqfullname}
\email{\lqemail @student.kit.edu}% chktex 1
% credit flo
\addtobeamertemplate{title page}{}{%
\scriptsize\color{gray}\begin{tikzpicture}[overlay, remember picture]
    \node[above left = 1em and 1em of current page.south east,align=center] (n0) {\LaTeX-Vorlage von\\\href{https://github.com/EagleoutIce}{Florian Sihler}};
    \node[above = 0cm of n0] (n1) {\fancyqr[size=1.5cm,tight]{https://eagleoutice.github.io/beamer-themes/\#colorful-dream}};
\end{tikzpicture}}% chktex 31

%%%%%%%%%%%%%%%%%%%%%%%%%%%%%%%%%%%%%%%%%%%%%%%%%%%%%%%%%%%%%%%%

\subtitle{Tutorium 3}
\date{02. Juni 2026}

\titleimagebox{2}{5.5cm}{
\begin{tikzpicture}[circuit logic IEC]
    \node (i0) {a};
    \node[below=of i0] (i1) {b};
    \node[below=of i1] (i2) {c};
    \node[below right=0 and 1 of i0, and gate] (a0) {};
    \node[below right=0 and 1 of a0, xor gate, xor gate IEC symbol={=}] (e0) {};
    
    \draw (i0.east)   -- ++(right:3mm) |- (a0.input 1);
    \draw (i1.east)   -- ++(right:3mm) |- (a0.input 2);
    \draw (a0.output) -- ++(right:3mm) |- (e0.input 1);
    \draw (i2.east)   -- ++(right:3mm) |- (e0.input 2);
    \draw (e0.output) -- ++(right:3mm);
\end{tikzpicture}
}

\outroimagetext{\begin{tabular}{c}
    Wtf ist das Würfelkalkül?
\end{tabular}}

%%%%%%%%%%%%%%%%%%%%%%%%%%%%%%%%%%%%%%%%%%%%%%%%%%%%%%%%%%%%%%%%

\begin{document}

\section[\(\overline{\wedge}\) hier, aber kennen Sie schon Näääänd?]{Blatt 03}

\subsection{Aufgabe 1}
\begin{frame}{Haben wir doch schon gemacht!}\onslide<+->
    \begin{enumerate}
        \item \(NAND_k\) ist dann \T{0}, wenn \onslide<+->alle Inputs \T{1} sind!
        \item<+-> \(NOT(x)=NAND_2(x,x)\)
        \item<+-> \(\overline\wedge \) ist nicht assoziativ!
        \item<+-> {\setlength{\abovedisplayskip}{0pt}\begin{align*}
            f(c,b,a) &= cba \vee \alert<.(1)>{(b\rightarrow c)}\\
            \only<+->{&= \alert<.(1)>{cba} \vee \alert<.>{\overline b \vee c}\\
            \only<+->{&= \overline b \vee \alert<.(1)>{c \vee \alert<.>{cba}}\\
            \only<+->{&= \overline b \vee \alert<.>{c}\\
            \only<+->{&= NAND_2(b,\overline c)}}}}
        \end{align*}}
    \end{enumerate}
\end{frame}

\subsection{Aufgabe 2}
\begin{frame}{Noch mehr von letztem Mal?!}\onslide<+->
    \begin{enumerate}
        \item Lesen! \onslide<+->\texttwemoji{eyes}
        \item<+-> \(K_1=%
        \only<.>{a\wedge\overline a\wedge b\wedge c}%
        \only<+>{\alert{a\wedge\overline a}\wedge b\wedge c}%
        \only<+>{0\wedge b\wedge c}%
        \only<+>{\alert{0}\wedge b\wedge c}%
        \only<+->{0}\) \\
        \onslide<+->\(K_2=%
        \only<.>{b\wedge a\wedge b\wedge c}%
        \only<+>{\alert{b}\wedge a\wedge \alert{b}\wedge c}%
        \only<+->{a\wedge b\wedge c}\) \\
        \onslide<+->\(K_3=%
        \only<.>{a\wedge 0\wedge b}%
        \only<+>{a\wedge \alert{0}\wedge b}%
        \only<+->{0}\) \\
        \onslide<+->\(K_4=%
        \only<.>{1\wedge a\wedge b}%
        \only<+>{\alert{1}\wedge a\wedge b}%
        \only<+->{a\wedge b}\)
        \item<+-> pro Literal: positiv, negativ, nicht im Term\\
        \onslide<+-> \T{1} ist die leere Konjunktion\\
        \onslide<+-> \T{0} muss noch dazu \onslide<+->\(\Rightarrow 3^2+1 = 10\) Produktterme
        \item<+-> \(a\wedge b,\quad a\wedge\overline b,\quad \overline a\wedge b,\quad \overline a\wedge \overline b\)
    \end{enumerate}
\end{frame}

\subsection{Aufgabe 3}
\begin{frame}{Von Treppen und Stufen\dots}\onslide<+->
    \begin{enumerate}
        \item<+-> {\setlength{\abovedisplayskip}{0pt}\begin{align*}
            f &= (a\rightarrow b)\wedge (b\rightarrow c)\\
            \only<+->{&= (\overline a\vee b)\wedge (\overline b \vee c)\\
            \only<+->{&= \overline a\overline b\vee \overline ac\vee b\overline b \vee bc\\
            \only<+->{&= \overline a\overline b \vee \overline ac \vee bc}}}
        \end{align*}}
        \item<+-> \(DNF = MINt(\dots)\)
        \item<+-> \(KNF = MAXt(\dots)\)
    \end{enumerate}
\end{frame}

\section[vorbereitet \texttwemoji{nerd_face}]{Tutorium}

\begin{frame}{Heute gibt es}
    \begin{itemize}
        \item Shannonscher Entwicklungssatz
        \item Min- und Maxterme (Wiederholung)
        \item Schaltfunktionen und Schaltnetze
    \end{itemize}
\end{frame}

\end{document}